\chapter{Architetture dei Modelli}

In questo capitolo vengono descritte le architetture dei modelli implementati. Tutti i modelli operano su immagini RGB normalizzate in $[-1, 1]$ e producono output nello stesso intervallo. Si distinguono due paradigmi: la \textbf{predizione diretta} dell'immagine pulita e il \textbf{residual learning}, dove il modello predice il rumore da sottrarre all'input.

% ============================================================
\section{UNet}
\label{sec:unet}

\textit{Sviluppato da: Giuseppe Bellamacina}

\medskip

La UNet \`e un'architettura encoder-decoder con skip connection, originariamente proposta per la segmentazione biomedica e adattata qui al restauro di immagini. Il modello apprende una mappatura diretta dall'immagine degradata a quella pulita.

\subsection{Architettura}

L'encoder \`e composto da 4 stadi di downsampling, ciascuno formato da un blocco \texttt{DoubleConv} (due convoluzioni $3\!\times\!3$ con BatchNorm e ReLU) seguito da MaxPool $2\!\times\!2$:
$$
\texttt{inc}(3 \to 64) \to \texttt{down}_1(64 \to 128) \to \texttt{down}_2(128 \to 256) \to \texttt{down}_3(256 \to 512) \to \texttt{down}_4(512 \to 1024)
$$

Il decoder ricostruisce la risoluzione spaziale tramite upsampling. Sono supportate due modalit\`a:
\begin{itemize}
    \item \textbf{Bilinear:} \texttt{nn.Upsample(scale\_factor=2, mode='bilinear')} seguito da \texttt{DoubleConv} con canali dimezzati.
    \item \textbf{Transposed Convolution:} \texttt{ConvTranspose2d(kernel\_size=2, stride=2)}.
\end{itemize}

Ad ogni livello del decoder, le feature dell'encoder vengono concatenate tramite skip connection dopo un eventuale padding per gestire discrepanze dimensionali. L'output finale \`e una convoluzione $1\!\times\!1$ che mappa da 64 a 3 canali, senza funzione di attivazione.

\subsection{Configurazione}

\begin{table}[H]
\centering
\begin{tabular}{ll}
\toprule
\textbf{Parametro} & \textbf{Valore} \\
\midrule
Canali di ingresso/uscita & 3 (RGB) \\
Feature iniziali & 64 \\
Modalit\`a upsampling & Bilinear / Transposed Conv \\
Paradigma & Predizione diretta \\
\bottomrule
\end{tabular}
\caption{Parametri architetturali della UNet.}
\end{table}

% ============================================================
\section{UNet Residual}
\label{sec:unet-residual}

\textit{Sviluppato da: Giuseppe Bellamacina}

\medskip

La UNet Residual condivide la stessa architettura encoder-decoder della UNet base, con una differenza fondamentale nel paradigma di apprendimento: il modello predice il \textbf{rumore} presente nell'immagine anzich\'e l'immagine pulita. L'output denoised si ottiene per sottrazione:
$$
\hat{x}_{\text{clean}} = x_{\text{noisy}} - f_\theta(x_{\text{noisy}})
$$

dove $f_\theta$ \`e la rete che predice la componente di rumore.

\subsection{Anti-checkerboard}

Nella variante con transposed convolution, \`e stato applicato un fix per gli artefatti a scacchiera (\textit{checkerboard artifacts}) tipici della deconvoluzione: il kernel della \texttt{ConvTranspose2d} \`e stato portato da $2\!\times\!2$ a $4\!\times\!4$ con \texttt{padding=1}, garantendo una copertura pi\`u uniforme dei pixel di output.

% ============================================================
\section{Attention UNet}
\label{sec:attention-unet}

\textit{Sviluppato da: Giuseppe Bellamacina}

\medskip

L'Attention UNet estende la UNet con meccanismi di \textit{attention gate} (AG) nelle skip connection, permettendo al decoder di concentrarsi sulle regioni pi\`u rilevanti delle feature dell'encoder.

\subsection{Attention Gate}

L'attention gate calcola una mappa di attenzione spaziale a partire dal segnale di gating $g$ (proveniente dal livello decoder inferiore) e dal segnale skip $x$ (proveniente dall'encoder):
$$
\alpha = \sigma\!\left(\psi\!\left(\text{ReLU}(W_g \cdot g + W_x \cdot x)\right)\right)
$$

dove $W_g$ e $W_x$ sono convoluzioni $1\!\times\!1$ con normalizzazione, $\psi$ \`e una convoluzione $1\!\times\!1$ che riduce a un singolo canale, e $\sigma$ \`e la funzione sigmoide. Le feature skip vengono modulate element-wise:
$$
\hat{x} = \alpha \odot x
$$

\subsection{Livelli di attenzione}

Il modello supporta due configurazioni:
\begin{itemize}
    \item \texttt{("bottleneck",)}: attention gate applicato solo al primo livello del decoder (collegamento con il bottleneck).
    \item \texttt{("all",)}: attention gate applicato a tutti i livelli del decoder.
\end{itemize}

\subsection{Normalizzazione}

Il modello supporta sia \texttt{BatchNorm2d} che \texttt{GroupNorm} (con $\min(8, C)$ gruppi), selezionabile tramite il parametro \texttt{norm}.

\subsection{Adaptive Blending}

L'Attention UNet supporta un meccanismo opzionale di blending adattivo basato sull'intensit\`a del rumore stimata. Se viene fornito un valore di \texttt{noise\_sigma} durante l'inferenza:
$$
\alpha = \min\!\left(\frac{\sigma_{\text{noise}}}{\sigma_{\text{trained}}},\; 1.0\right), \qquad
\text{output} = \alpha \cdot x_{\text{input}} + (1 - \alpha) \cdot x_{\text{denoised}}
$$

Questo permette di moderare l'intervento del modello quando il rumore \`e inferiore a quello su cui \`e stato addestrato.

% ============================================================
\section{DnCNN}
\label{sec:dncnn}

\textit{Sviluppato da: Salvo Iurato}

\medskip

Il DnCNN (\textit{Denoising Convolutional Neural Network}) \`e un'architettura puramente sequenziale che sfrutta il residual learning. A differenza delle architetture encoder-decoder, non utilizza alcuna forma di downsampling o skip connection.

\subsection{Architettura}

La rete \`e composta da 20 layer convoluzionali:
\begin{enumerate}
    \item \textbf{Primo layer:} \texttt{Conv2d(3$\to$64, 3$\times$3)} con bias + ReLU (senza BatchNorm).
    \item \textbf{Layer 2--19:} \texttt{Conv2d(64$\to$64, 3$\times$3)} senza bias + BatchNorm + ReLU.
    \item \textbf{Ultimo layer:} \texttt{Conv2d(64$\to$3, 3$\times$3)} con bias (senza BatchNorm n\'e attivazione).
\end{enumerate}

Il rumore predetto viene sottratto dall'input:
$$
\hat{x}_{\text{clean}} = x_{\text{noisy}} - f_\theta(x_{\text{noisy}})
$$

In fase di inferenza, l'output viene clampato nell'intervallo $[-1, 1]$. I pesi sono inizializzati con Kaiming normal per le convoluzioni e costanti (1 / 0) per la BatchNorm.

% ============================================================
\section{Denoising Autoencoder}
\label{sec:denoising-autoencoder}

\textit{Sviluppato da: Salvo Iurato}

\medskip

Il Denoising Autoencoder \`e un'architettura encoder-decoder \textbf{senza skip connection}, progettata come baseline per il confronto con le architetture pi\`u avanzate.

\subsection{Architettura}

L'\textbf{encoder} esegue 4 stadi di downsampling tramite convoluzioni strided ($3\!\times\!3$, \texttt{stride=2}), ciascuno seguito da BatchNorm e LeakyReLU(0.2):
$$
3 \xrightarrow{128\to 64} 32 \xrightarrow{64\to 32} 64 \xrightarrow{32\to 16} 128 \xrightarrow{16\to 8} 256
$$

Il \textbf{decoder} ricostruisce la risoluzione tramite \texttt{Upsample(bilinear)} + \texttt{Conv2d(3$\times$3)} per ogni stadio, con BatchNorm e LeakyReLU, fino a un layer finale con attivazione Tanh che produce l'output in $[-1, 1]$.

L'assenza di skip connection costringe il bottleneck a contenere tutta l'informazione necessaria per la ricostruzione, rendendo il modello pi\`u ``lossy'' ma anche pi\`u semplice.

% ============================================================
\section{RIDNet}
\label{sec:ridnet}

\textit{Sviluppato da: Salvo Iurato}

\medskip

RIDNet (\textit{Real Image Denoising Network}) \`e un modello che combina convoluzioni dilatate e attention a livello di canale per catturare informazioni multi-scala senza ricorrere a operazioni di pooling.

\subsection{Enhancement Attention Block (EAB)}

L'unit\`a fondamentale di RIDNet \`e l'\textit{Enhancement Attention Block}, composto da:
\begin{enumerate}
    \item Tre convoluzioni con dilatazione crescente ($d = 1, 2, 3$), ciascuna seguita da BatchNorm e ReLU. Questo expande il campo recettivo senza ridurre la risoluzione spaziale.
    \item Un modulo \texttt{ChannelAttention} di tipo Squeeze-and-Excitation: Global Average Pooling $\to$ Conv $1\!\times\!1$ ($C \to C/r$) $\to$ ReLU $\to$ Conv $1\!\times\!1$ ($C/r \to C$) $\to$ Sigmoid, con fattore di riduzione $r = 16$.
    \item Una connessione residuale locale con fattore di scala:
\end{enumerate}
$$
\text{EAB}(x) = x + s \cdot \text{CA}(\text{dilated\_branch}(x))
$$

dove $s$ \`e il \texttt{residual\_scale} (default 1.0, negli esperimenti 0.1).

\subsection{Architettura completa}

\begin{enumerate}
    \item \textbf{Estrazione feature:} \texttt{Conv2d(3$\to$64, 3$\times$3)} + ReLU.
    \item \textbf{Enhancement:} Sequenza di $N$ EAB (tipicamente $N = 8$).
    \item \textbf{Ricostruzione:} \texttt{Conv2d(64$\to$3, 3$\times$3)} che predice il rumore.
    \item \textbf{Residual globale:} $\hat{x} = x - f_\theta(x)$.
\end{enumerate}

In inferenza l'output viene clampato a $[-1, 1]$.

% ============================================================
\section{Pix2Pix}
\label{sec:pix2pix}

\textit{Sviluppato da: Daniele Barbagallo}

\medskip

Pix2Pix \`e un framework di \textit{image-to-image translation} basato su Generative Adversarial Network (GAN) condizionale. Il sistema \`e composto da un \textbf{generatore} e un \textbf{discriminatore} addestrati in modo avversariale.

\subsection{Generatore}

Il generatore ha un'architettura UNet-like con 7 livelli di downsampling, progettata per input $256\!\times\!256$:

\textbf{Encoder} (6 layer + bottleneck):
$$
\texttt{enc}_1(3 \to 64) \to \texttt{enc}_2(64 \to 128) \to \cdots \to \texttt{enc}_6(512 \to 512) \to \texttt{bottleneck}(512 \to 512)
$$

Ogni layer usa \texttt{Conv2d(4$\times$4, stride=2)} + BatchNorm + LeakyReLU(0.2). Il primo layer non ha BatchNorm.

\textbf{Decoder} (6 layer + finale): \texttt{ConvTranspose2d(4$\times$4, stride=2)} + BatchNorm + ReLU, con Dropout(0.5) nei primi 3 layer. Skip connection via concatenazione con il corrispondente encoder. Il layer finale produce 3 canali con attivazione Tanh (o senza, per residual learning).

\subsection{PatchGAN Discriminator}

Il discriminatore opera su \textit{patch} locali dell'immagine, classificando regioni $70\!\times\!70$ come reali o generate. L'input \`e la concatenazione dell'immagine condizionale (degradata) e dell'immagine target/generata (6 canali totali).

L'architettura consiste in 4 blocchi convoluzionali ($4\!\times\!4$, stride progressivamente ridotto) che riducono le feature da 64 a 512 canali, seguiti da una convoluzione finale $4\!\times\!4$ che produce un output a singolo canale (logits, senza sigmoide).

\subsection{Modalit\`a di apprendimento}

Nell'esperimento condotto, il generatore opera in \textbf{residual learning}: predice $\hat{n} = G(x_{\text{noisy}})$ (il rumore stimato) e l'immagine pulita si ottiene come:
$$
\hat{x} = \text{clamp}(x_{\text{noisy}} - \hat{n},\; -1,\; 1)
$$

Il rumore predetto viene inoltre clampato a $\pm 0.6$ per stabilit\`a.

% ============================================================
\section{NAFNet}
\label{sec:nafnet}

\textit{Sviluppato da: Mattia Campanella}

\medskip

% TODO: Sezione da completare dopo il merge del branch di Mattia
\textit{Sezione da completare --- codice non ancora presente sul branch main.}

% ============================================================
\section{GAN}
\label{sec:gan}

\textit{Sviluppato da: Daniele Barbagallo}

\medskip

% TODO: Sezione da completare dopo il merge del branch di Daniele (GAN aggiuntiva rispetto a Pix2Pix)
\textit{Sezione da completare --- codice non ancora presente sul branch main.}

% ============================================================
\section{Riepilogo delle architetture}

\begin{table}[H]
\centering
\small
\begin{tabular}{lcccc}
\toprule
\textbf{Modello} & \textbf{Tipo} & \textbf{Residual} & \textbf{Attention} & \textbf{Attivazione output} \\
\midrule
UNet & Enc-Dec + skip & No & No & Nessuna \\
UNet Residual & Enc-Dec + skip & S\`i & No & Nessuna \\
Attention UNet & Enc-Dec + skip & No & AG (bottleneck/all) & Nessuna \\
DnCNN & Sequenziale (20 layer) & S\`i & No & Clamp $[-1,1]$ \\
Denoising AE & Enc-Dec (no skip) & No & No & Tanh \\
RIDNet & Feature + EAB & S\`i & Channel (SE) & Clamp $[-1,1]$ \\
Pix2Pix Gen. & UNet-like (7 livelli) & Configurabile & No & Tanh / Nessuna \\
\bottomrule
\end{tabular}
\caption{Confronto delle architetture implementate.}
\label{tab:architetture}
\end{table}
